\section{A reporting protocol for cross-lingual document VLM evaluation}
\label{sec:protocol}

The findings in \S\S\ref{sec:measurement}--\ref{sec:sft} are properties of \emph{strict-EM applied to cross-lingual generation}, not of MTVQA, XFUND, or the three architectures we tested. Strict-EM-only reporting is structurally biased: it conflates capability, output-style convention, and per-language gold-annotation granularity, and it can reverse the cross-VLM ranking on identical inputs. We propose a concrete reporting protocol that any cross-lingual document VLM evaluation can adopt at single-digit dollar cost.

\paragraph{Nine required items (R1--R9).}
A multilingual document VLM evaluation reporting a cross-lingual finding must supply: \textbf{(R1)} per-language strict-EM with $n$ and Wilson $95\%$ CI; \textbf{(R2)} per-language semantic-EM-C and -C+P with $n$ and $95\%$ CI (vision-grounded LLM-judge adjudication of (image, question, gold, prediction) tuples); \textbf{(R3)} judge specification --- model name and version, decoding settings, verbatim adjudication prompt, verdict-parsing logic; \textbf{(R4)} judge agreement statistic --- human on $\geq 50$ stratified rows ($\kappa \geq 0.6$ or 3-way $\geq 80\%$) and/or cross-judge agreement with a second judge from a different model family at the same threshold; \textbf{(R5)} retry-exhaustion accounting per language plus sensitivity to drop/generous/conservative coercions (cf.\ Appendix~\ref{app:dropped-sensitivity}); \textbf{(R6)} the \emph{output-style penalty} (per-language strict$-$sem-C, reported as a diagnostic); \textbf{(R7)} cross-VLM comparisons on the \emph{intersection} of sample\_ids that received clean verdicts on every VLM; \textbf{(R8)} per-language $\Delta$strict-EM \emph{and} $\Delta$semantic-EM-C for every intervention claim --- the pattern (\emph{strict lifts, semantic flat} or \emph{strict lifts, semantic drops}) is the diagnostic for normalisation vs capability; \textbf{(R9)} public release of (image\_uid, gold, prediction, verdict, reason) tuples plus verdict-parsing code and a script that reproduces every aggregate cell. A worked example applying R1--R9 to a hypothetical ``$+5.0$ pp Arabic strict-EM lift'' claim is in Appendix~\ref{app:protocol-example}.

\paragraph{Cost and scope.}
At standard 2026 proprietary-judge rates (\$0.014 per tuple via Gemini-2.5-Flash on Vertex AI, including a 6-attempt retry budget), the full MTVQA val ($n=2203$) costs \$31; a per-language $n=300$ sample costs \$4; a 358-row held-out costs \$5. The open-source variant (Qwen2.5-7B-Instruct in vision mode, single L4 GPU, $\sim 2$ h for 358 rows) eliminates monetary cost entirely; cross-judge agreement is the validation criterion, not judge identity. Strict-EM is not abolished: it remains useful as a backward-comparable column, the primary metric for character-precision tasks (IDs, dates, totals), and a conservative lower-bound capability proxy. The protocol \emph{adds} semantic-EM, the output-style penalty, and the joint-set discipline; it \emph{subtracts} no existing metric.

The recommendation: at the next round of MTVQA, DocVQA-multilingual, InfographicVQA-multilingual, and similar benchmark refreshes, R1--R9 should be the submission template. Vision-grounded LLM-as-judge at sufficient reliability has only been available since late 2024; the cost-benefit analysis is the most acute on non-Latin script generation where gold-string surface form varies enough to invert leaderboard rankings, as we have shown.
